\section{ISO 21448: Sigurnost namijenjene funkcionalnosti (SOTIF)}
\label{sec:sotif}

\subsection{Motivacija i definicija}

Norma ISO~21448 \cite{iso21448}, poznata pod akronimom SOTIF (engl.\ \textit{Safety of the Intended Functionality}), objavljena je 2022.\ godine kao odgovor na spoznaju da funkcijska sigurnost prema ISO~26262 ne pokriva sve izvore opasnosti relevantne za sustave automatizirane vožnje. Norma je evoluirala iz dokumenta ISO/PAS~21448, objavljenog 2019.\ godine, što odražava brzo sazrijevanje područja \cite{patel2025sotifslr}. Dok ISO~26262 adresira rizike uzrokovane \textit{kvarovima} sustava, SOTIF se bavi opasnostima koje nastaju kada sustav radi \textit{točno onako kako je projektiran}, ali njegove performanse nisu dovoljne za sigurno funkcioniranje u svim operativnim uvjetima \cite{patel2025sotifslr}.

SOTIF definira svoj cilj kao \textit{odsutnost neprihvatljivog rizika uzrokovanog opasnostima koje proizlaze iz funkcijskih insuficijencija ili razumno predvidive zlouporabe sustava} \cite{iso21448}. Ova definicija obuhvaća dva ključna izvora opasnosti:

\begin{itemize}
  \item \textbf{Funkcionalne insuficijencije} — ograničenja u performansama senzora, algoritama percepcije ili strategija donošenja odluka. Primjer: kamera koja ne prepoznaje pješaka u uvjetima jake kiše ili sustav koji pogrešno interpretira odbljesak sunca kao prepreku;
  \item \textbf{Razumno predvidiva zlouporaba} — situacije u kojima korisnik sustav koristi na način koji nije predviđen, ali je očekivan. Primjer: prekomjerno oslanjanje vozača na sustav razine SAE~2 uz nedovoljno praćenje prometne situacije.
\end{itemize}

\subsection{Područje primjene}

ISO~21448 primjenjuje se na sustave razina automatizacije SAE~1 do SAE~5 \cite{saej3016}. Norma \textit{ne} pokriva \cite{iso21448}:
\begin{itemize}
  \item Kvarove hardvera i softvera (to pokriva ISO~26262);
  \item Kibernetičku sigurnost (pokrivena normom ISO/SAE~21434);
  \item Opasnosti inherentne samoj tehnologiji (npr.\ lasersko zračenje LiDAR-a).
\end{itemize}

Primjena SOTIF-a posebno je relevantna za sustave koji koriste napredne senzorske tehnologije (kamere, LiDAR, radar), algoritme strojnog učenja za percepciju i sustave fuzije senzora — sve komponente kod kojih postoji inherentna nesigurnost u interpretaciji okoline \cite{patel2025sotifslr}. Kako je prikazano u poglavlju~\ref{sec:percepcija}, upravo perceptivni lanac predstavlja glavni izvor funkcionalnih insuficijencija.

\subsection{Operativna domena (ODD)}

Temeljni koncept za primjenu SOTIF-a je operativna domena (engl.\ \textit{Operational Design Domain}, ODD) — skup uvjeta pod kojima je sustav projektiran da radi sigurno, uključujući tipove cesta, geografska područja, raspon brzina, vremenske i svjetlosne uvjete. Scenariji unutar ODD-a klasificiraju se prema potencijalu da uzrokuju opasno ponašanje, a izlazak izvan ODD-a (ili rad na njegovim rubovima) tipičan je izvor SOTIF rizika \cite{patel2025sotifslr}. Definiranje i nadzor ODD-a stoga su preduvjet za smislenu procjenu sigurnosti.

\subsection{Okidajući uvjeti}

Ključni koncept u SOTIF-u su \textit{okidajući uvjeti} (engl.\ \textit{triggering conditions}) — specifične okolišne ili situacijske okolnosti koje mogu uzrokovati da sustav koji radi prema specifikaciji ipak proizvede opasan ishod. Okidajući uvjeti mogu biti \cite{iso21448, patel2025sotifslr}:

\begin{itemize}
  \item \textbf{Okolišni} — nepovoljni vremenski uvjeti (kiša, magla, snijeg), nedostatno ili prekomjerno osvjetljenje, refleksije;
  \item \textbf{Scenarijski} — neočekivani obrasci ponašanja drugih sudionika u prometu, netipične konfiguracije prometnica;
  \item \textbf{Tehnološki} — ograničenja rezolucije senzora, degradacija performansi uslijed zagađenja senzora (npr.\ blato na kameri).
\end{itemize}

Veza s percepcijom ovdje je izravna: kako pokazuje Tablica~\ref{tab:senzori-uvjeti} u poglavlju~\ref{sec:percepcija}, magla, jaka kiša i zasljepljujuće sunce mjerljivo degradiraju kamere i (u manjoj mjeri) LiDAR. Identifikacija i klasifikacija okidajućih uvjeta temeljna je aktivnost SOTIF analize jer omogućuje sustavno smanjivanje prostora \textit{nepoznatih opasnih scenarija}.

\subsection{Klasifikacija scenarija}

SOTIF uvodi klasifikaciju operativnih scenarija u četiri područja, prema dvjema osima — jesu li scenariji poznati i jesu li opasni \cite{patel2025sotifslr, iso21448}:

\begin{enumerate}
  \item \textbf{Područje 1 — poznati, neopasni scenariji};
  \item \textbf{Područje 2 — poznati, opasni scenariji};
  \item \textbf{Područje 3 — nepoznati, opasni scenariji};
  \item \textbf{Područje 4 — nepoznati, neopasni scenariji}.
\end{enumerate}

Cilj SOTIF procesa je analizom namijenjene funkcionalnosti, modifikacijom sustava te verifikacijom i validacijom procijeniti i smanjiti rizik opasnog ponašanja u Području~2 i Području~3 \cite{patel2025sotifslr}. Poznati opasni scenariji (Područje~2) zahtijevaju implementaciju mjera koje ih prebacuju u sigurno područje, dok je za nepoznate opasne scenarije (Područje~3) — najizazovniju kategoriju — potrebno sustavnim testiranjem i analizom prevoditi ih u poznate.

\subsection{Mjere za osiguravanje SOTIF-a}

SOTIF mjere uključuju modifikaciju funkcionalnosti sustava ili zahtjeva na performanse senzora, implementaciju redundancije i fail-safe mehanizama te temeljito testiranje i validaciju \cite{patel2025sotifslr}. Konkretno, modifikacije sustava obuhvaćaju poboljšanje performansi ili točnosti senzora; redundancija podrazumijeva dodatne senzore koji održavaju namijenjenu funkcionalnost u slučaju zatajenja; a fail-safe mehanizmi (npr.\ sustav za nužno kočenje) automatski se aktiviraju u opasnim situacijama. Te se mjere međusobno nadopunjuju kako bi se rizik opasnog ponašanja sveo na prihvatljivu razinu.

\subsection{Iterativni pristup verifikaciji i validaciji}

Za razliku od V-modela u ISO~26262, SOTIF koristi iterativni pristup koji ne slijedi strogo linearnu progresiju. Proces uključuje \cite{iso21448, patel2025sotifslr}:

\begin{enumerate}
  \item \textbf{Specifikaciju i analizu funkcionalnosti} — definicija namijenjenog ponašanja sustava i identifikacija potencijalnih funkcionalnih insuficijencija;
  \item \textbf{Identifikaciju i procjenu opasnih scenarija} — analiza mogućih kombinacija okidajućih uvjeta i funkcionalnih insuficijencija;
  \item \textbf{Definiciju strategija za smanjenje rizika} — modifikacije dizajna, ograničenja operativne domene (ODD), upozorenja korisniku;
  \item \textbf{Verifikaciju i validaciju} — potvrda da implementirane mjere učinkovito smanjuju rezidualni rizik na prihvatljivu razinu, kroz testiranje u zatvorenoj petlji sa softverom (engl.\ \textit{software-in-the-loop}), hardverom (engl.\ \textit{hardware-in-the-loop}) i na razini vozila \cite{patel2025sotifslr};
  \item \textbf{Procjenu rezidualnog rizika} — kvantitativna ili kvalitativna ocjena preostalog rizika nakon primjene svih mjera.
\end{enumerate}

Ovaj iterativni ciklus ponavlja se sve dok rezidualni rizik ne zadovolji definirane kriterije prihvatljivosti.

\subsection{Metode analize opasnosti i scenarijsko testiranje}

Sustavni pregled literature o SOTIF-u \cite{patel2025sotifslr} prepoznaje više skupina metoda za identifikaciju opasnosti i procjenu rizika prilagođenih ADS-u:
\begin{itemize}
  \item \textbf{STPA} (engl.\ \textit{System-Theoretic Process Analysis}) — identificira nesigurne upravljačke akcije i uzročne čimbenike te ih povezuje sa sigurnosnim zahtjevima izvedenima iz ograničenja ML komponenata;
  \item \textbf{Bayesove mreže} — modeliraju i kvantificiraju nesigurnosti u performansama sustava te računaju stope opasnosti za različite konfiguracije;
  \item \textbf{Kvantitativna norma rizika} (engl.\ \textit{Quantitative Risk Norm}, QRN) — prilagođava HARA-u za ADS definiranjem prihvatljivih frekvencija incidenata po razredima ozbiljnosti;
  \item \textbf{FTA i ETA} (analiza stabla kvarova i stabla događaja) — proširene za primjenu na SOTIF;
  \item \textbf{Detekcija scenarija izvan distribucije} (engl.\ \textit{out-of-distribution}) — razlikuje viđene od neviđenih podataka i okida rezervne mehanizme.
\end{itemize}
Scenarijsko testiranje pritom kombinira varijable iz više izvora — slučajeve uporabe, okidajuće uvjete, predvidivu zlouporabu i ograničenja senzora — uz naglasak na simulacijski vođeno testiranje i povratne petlje vođene dokazima \cite{patel2025sotifslr}.

\subsection{Sigurnost percepcije temeljene na strojnom učenju}

Poseban izazov za SOTIF predstavljaju komponente percepcije temeljene na strojnom učenju. Njihovi netočni izlazi mogu ugroziti sigurnost, a izvori problema uključuju pristranosti u podacima za učenje, nepotpune skupove podataka te slab nastup na rijetkim, rubnim slučajevima \cite{patel2025sotifslr}. Literatura predlaže izvođenje sigurnosnih zahtjeva za ML-percepciju (npr.\ STPA-pristupom), okvire za osiguranje strojnog učenja te detekciju scenarija izvan distribucije kao mehanizam za upravljanje nesigurnošću \cite{patel2025sotifslr}.

\subsection{Zona prihvatljivog rezidualnog rizika i otvoreni izazovi}

SOTIF prepoznaje da potpuna eliminacija rizika nije moguća — umjesto toga, cilj je postići \textit{prihvatljiv rezidualni rizik}. Ovo zahtijeva definiciju kvantitativnih ili kvalitativnih kriterija prihvatljivosti, što predstavlja jedan od najvećih izazova u primjeni norme. Sustavni pregled literature kao ključne istraživačke praznine ističe nedostatak \textit{jedinstvenih kriterija evaluacije} SOTIF mjera te ograničeno razumijevanje kako u SOTIF metodologije integrirati \textit{ljudske čimbenike} \cite{patel2025sotifslr}. Trenutna istraživanja stoga su usmjerena na razvoj kvantitativnih metoda validacije koje bi omogućile objektivnu i usporedivu procjenu dosegnutog stupnja sigurnosti.
